\documentclass[a4paper]{article}

\usepackage{a4wide}
\usepackage[utf8]{inputenc}
\usepackage[ngerman]{babel}
\usepackage{fancyhdr}
\usepackage{graphicx}
\usepackage[parfill]{parskip}
\usepackage{color}

% fuer "deko-elemente"
\pagestyle{fancy}

% we needs moar platz
\topmargin -5mm
\headheight 20mm

% include logo in document head
\chead{\includegraphics[width=50mm]{logo.pdf}}

% seitenzahl entfernen
\cfoot{}

% footer content
\lfoot{\small{AlphaLabs}}
\rfoot{\tiny{Alexander Drangmeister, Dennis Esders, Philipp Heisig, Robert Baruck, Stephan Wilde, Kilian Költzsch}}

% headrule
\renewcommand\headrule{
\vspace{3mm}
\footrule
}

% footrule
\renewcommand\footnoterule{\rule{\linewidth}{0pt}}

% change standard font
\usepackage[T1]{fontenc}
\newcommand{\changefont}[3]{
\fontfamily{#1} \fontseries{#2} \fontshape{#3} \selectfont}

% Tagesordnungspunkt Item
\newcommand{\TOP}[1]{\item \textbf{#1}\par}



\begin{document}
\changefont{cmss}{m}{n} %computer modern sans serif

\begin{center}
\textbf{\Large Protokoll zum Gruppentreffen vom 16. April 2014}
\end{center}
\vspace{3mm}

\textbf{Beginn:} 16:00 Uhr, \textbf{Ende:} 17:00 Uhr\\
\textbf{Anwesende:} Alex, Dennis, Philipp, Robert, Stephan, Kilian

\vspace{5mm}

\begin{enumerate}

\TOP{Terminfindung}
Wir teilen Themengebiete auf die Gruppenmitglieder auf und treffen uns zum 23.04. erneut um die Ergebnisse zusammenzutragen und die Abgabe des nächsten Teilziels vorzubereiten.

\TOP{Spielname}
Der Spielname steht jetzt als \textit{The Circle of Waste} fest. So lehnen wir den Namen direkt an die Idee des Spiels, nämlich die Darstellung des Müllkreislaufs, an.

\TOP{Zielgruppe}
Um die Zielgruppe genauer zu definieren haben wir uns auf 12-jährige Schüler geeinigt. Diese bieten sich ideal an, da die Thematik für sie im Vergleich zu jüngeren Schülern relevanter ist und sich so unterrichtsbegleitend zu Sach- und Naturkunde gut angliedert.

\TOP{Zeitplan}
Zum nächsten Treffen soll ein grober Zeitplan zur besseren Planung des Projekts geschaffen werden.

\TOP{Mockups}
Ebenfalls werden erste Mockups gestaltet, um das Spiel vor der eigentlichen Implementation besser zu visualisieren.

\TOP{Minispielideen}
Da das Spiel aus diversen Minispielen bestehen soll, werden hier auch bis zum nächsten Treffen konkretere Vorschläge entstehen.

\end{enumerate}


% das ist herrlich kompliziert, aber es klappt! das ist alles was zaehlt! :D
\definecolor{dgreen}{rgb}{0.549,0.776,0.247}
\makeatletter
\def\footrule{{
  \vskip-\footruleskip\vskip-\footrulewidth
  \color{\footrulecolor}
  \hrule\@width\headwidth\@height
  \footrulewidth\vskip\footruleskip
}}
\makeatother
\renewcommand{\footrulewidth}{3pt}
\newcommand{\footrulecolor}{dgreen}

\end{document}
